\section{Threats to Validity}
\label{sec:threats}

\paragraph{Internal threats.} Four things could make our comparison wrong.
First, we wrote the baseline ourselves, so it could be weaker than it should
be. Second, the live runner and the resource pool are small stand-ins that we
wrote for this evaluation, not the team's engine. Third, the same author wrote
the value generator and the schema checker of \rqref{2}, so both could share
one misreading of a keyword. Fourth, generation and the APIs are random. Our
guards:
\begin{itemize}
  \item the baseline follows the schema \texttt{type} and \texttt{enum},
    fills the same parameters as \tool, and goes through the same runner with
    the same budget; it is not tuned to fail;
  \item the runner and pool are the same for all four configurations, so the
    differences between configurations come from \tool;
  \item the checker parses dates and UUIDs with Python's standard library, not
    with code from \tool, and \rqref{2} counts only keywords that the checker
    understands;
  \item every configuration ran \NReps times with fixed seeds, each time on a
    freshly started API, and Table~\ref{tab:live} reports the standard
    deviation.
\end{itemize}
What remains: a real engine chooses requests from feedback, and may narrow or
widen the gaps we measured. Our 5XX signature masks quoted text in error
messages; two faults that differ only in a quoted name count as one, so
\Faults may be smaller than the number of distinct messages. In a trial run
we read the signatures of the operation with the most 5XX responses and found
no duplicates caused by echoed input.

\paragraph{External threats.} We ran \tool live on two APIs, both written in
Java with Spring, out of the five known REST League APIs, and the competition
adds five APIs that nobody has seen. All \NSpecs specifications are
OpenAPI~3.0 files in JSON. \tool reads YAML only when the PyYAML package is
present and does not read Swagger~2.0 at all, so an unseen API with such a
specification would fail at Step~1. None of the specifications had a
recursive schema. We used one machine and one version of each tool, listed in
Section~\ref{sec:setup}.

\paragraph{Construct threats.} We measure what REST League scores, but the
scores are proxies. An operation counts as reached after one 2XX response,
even when that response did little, such as a list request on an empty
collection. A unique 5XX response is not always a bug that a developer would
fix: many of the 5XX responses we saw on pet-clinic are wrong parameter types
that the API reports as 500 instead of 400. They count in the competition, so
we count them, but we do not call them bugs. Schema conformance in \rqref{2}
means only that a value fits the schema; the API can still reject it, which
is why \rqref{3} measures acceptance on running APIs. We did not measure code
coverage, the third REST League metric, because our runner talks to the APIs
directly instead of through RESTgym's coverage collection.
